% ============================================================
%  sec-3-operator-migration.tex  —  4.3 Operator Migration
% ============================================================

\section{Critical Operator Migration}
\label{sec:operator-migration}

The two most technically demanding components of the migration are
the Mamba \acrshort{ssm} selective scan kernel and the \acrshort{moe}
token dispatch kernel. Both are implemented as custom \acrshort{cuda}
extensions in the original pipeline and have no pre-built \acrshort{cann}
equivalent. This section documents the design and strategy for the
Ascend-compatible replacements provided in \filepath{ops/mamba\_ascend.py}
and \filepath{ops/moe\_routing\_ascend.py}.

% ── 4.3.1 ────────────────────────────────────────────────────
\subsection{Mamba SSM Selective Scan Operator}
\label{sec:mamba-operator}

\subsubsection{Original CUDA Implementation}

In the Nemotron 3 Nano hybrid architecture, Mamba layers implement
the selective state space recurrence via a hand-written \acrshort{cuda}
kernel (\texttt{selective\_scan\_cuda}). This kernel implements a parallel
prefix scan algorithm that computes the full recurrence across all
$L$ timesteps in $O(\log L)$ depth, exploiting the associativity
of the state update to parallelise across the sequence dimension.

\subsubsection{CANN Incompatibility}

The \texttt{selective\_scan\_cuda} kernel cannot execute on Ascend
because it is compiled as a \acrshort{cuda} \acrshort{ptx} binary.
Furthermore, the parallel prefix scan algorithm relies on CUDA warp
shuffle instructions which have no direct equivalent in CANN's DaVinci
instruction set.

\subsubsection{PyTorch Fallback Implementation}

The replacement in \filepath{ops/mamba\_ascend.py} provides two
implementations: a sequential reference and an optimised chunked
version. The sequential version implements the recurrence in a Python
loop over timesteps, while the chunked implementation reduces Python
overhead by processing the sequence in blocks of 256 tokens. The
actual implementation code is provided in Appendix A.

The chunked version is selected automatically when the environment
variable \texttt{USE\_CHUNKED\_MAMBA=1} is set, which is the default
in the Ascend cluster profile.

\subsubsection{Kernel Injection Mechanism}

The replacement is injected at process startup in
\filepath{train\_sft\_ascend.py} by reassigning module-level function
pointers in the \texttt{mamba\_ssm} package, ensuring all subsequent
calls invoke the Ascend-compatible fallback.

\subsubsection{Correctness Verification}

The built-in self-test in \filepath{ops/mamba\_ascend.py} validates
output shape correctness and numerical equivalence between the
implementations. \tabref{mamba-perf} summarises the characteristics of
all three implementations.

\begin{table}[htbp]
  \centering
  \caption{Mamba SSM implementation comparison}
  \label{tab:mamba-perf}
  \begin{tabularx}{\textwidth}{lXXX}
    \toprule
    \tableheader{Implementation & Correctness & Relative Speed & Device} \\
    \midrule
    \texttt{selective\_scan\_cuda}
      & Reference (ground truth)
      & 1.0$\times$ (baseline)
      & CUDA only \\
    \rowcolor{TableRowOdd}
    \texttt{selective\_scan\_ascend}
      & Matches CUDA (atol $<$ 1e-4)
      & $\approx$0.2$\times$
      & Any (CANN, CPU) \\
    \texttt{selective\_scan\_chunked}
      & Matches sequential (atol $<$ 1e-5)
      & $\approx$0.5$\times$
      & Any (CANN, CPU) \\
    \bottomrule
  \end{tabularx}
\end{table}

% ── 4.3.2 ────────────────────────────────────────────────────
\subsection{MoE Token Dispatch Operator}
\label{sec:moe-operator}

\subsubsection{Original CUDA Implementation}

Nemotron 3 Nano's sparse \acrshort{moe} layers activate 6 out of
128 experts per token using fused \acrshort{cuda} kernels for
routing weight computation and token dispatch (gathering/scattering).

\subsubsection{PyTorch Fallback Implementation}

The replacement in \filepath{ops/moe\_routing\_ascend.py} decomposes
the dispatch into a routing weight computation and an expert execution
loop. The implementation uses standard PyTorch operations to execute
each expert \acrshort{ffn} on its assigned token subset. The complete
source code is available in Appendix A.

\subsubsection{Load Balancing Loss}

The load balancing auxiliary loss from Switch Transformer prevents
expert collapse during fine-tuning by penalising uneven token
distribution across experts.

\subsubsection{Megatron-Bridge Dispatcher Patch}

The MoE dispatch replacement is injected into Megatron-Bridge's
\texttt{MoETokenDispatcher} class via monkey-patching at process
startup.

% ── 4.3.3 ────────────────────────────────────────────────────
\subsection{Autocast and Mixed-Precision Patch}
\label{subsec:autocast-patch}

Megatron-Bridge uses \texttt{torch.cuda.amp.autocast} for mixed
precision training. On Ascend, this is replaced with an
NPU-compatible wrapper in \filepath{train\_sft\_ascend.py} that
redirects calls to \texttt{torch.npu.amp.autocast()}, ensuring
correct FP16 acceleration on DaVinci Vector Units.
